\subsection{Data}\label{subsec:data-sources-and-data-integration}

The Data phase in Figure~\ref{fig:pipeline} turns heterogeneous authorisation exports into a shared and auditable data layer.

We use four authorisation sources: the European Union, Spain, France, and Portugal~\cite{ec-pesticides-database, mapa-regfiweb,anses-ephy,dgav-sifito}.
Each source has a different structure. We therefore create one parser per source and map the available fields to a shared schema.

The shared schema includes crop, pathogen, product, microbial active substance, product function, authorisation status, dates, source, and review information.
Fields that are not present in a source remain empty.
The pipeline does not infer them.

The sources do not provide the same type of evidence.
The EU data contain active substances and authorisation information, but no complete crop--pathogen use relation~\cite{eu-dataset-snapshot}.
The Spanish files provide product, crop, and substance catalogues, but no registration key that safely links products to uses~\cite{spain-dataset-snapshot}.
The French E-Phy export provides product-use rows, but not the scientific crop--pathogen pairs required for modelling ~\cite{france-dataset-snapshot}.
The Portuguese exports provide the most complete fields and supply the current joined modelling data~\cite{dgav-sifito}.

During import, the pipeline checks required columns and keeps optional fields when available.
It normalises accents, punctuation, and strain wording.
It also expands Portuguese rows that contain several crop or target names.
Ambiguous pairs are kept for review.
EPPO names and codes support the reference checks~\cite{eppo-global-database}.
